\section{De palabra de moda a punto de inflexión industrial: por qué 2025 se convirtió en el año de los Agentes}

\subsection{Tres realidades que sostienen el auge}
El juicio que dio el moderador en la apertura fue directo: 2025 no es la primera vez que aparece el concepto de Agent, pero sí es la primera vez que coinciden simultáneamente el capital, el producto y la investigación. La duplicación del tamaño del mercado, el avance hacia la salida a bolsa de las empresas líderes y el estallido paralelo de artículos académicos y sistemas de código abierto son señales que hicieron que el Agent pasara de ser un «concepto de demo» a una vía en la que apuestan tanto empresas como instituciones de investigación.

\begin{importantbox}{Por qué esta vez no es ruido pasajero}
Lo distinto de esta ola de auge de los Agentes es que ya no depende solo de una tabla de clasificación de modelos ni del lanzamiento de un solo producto, sino que apareció una cadena completa:
\begin{itemize}
  \item Aguas arriba hay LLM más fuertes, reasoning models y modelos de computer use;
  \item Aguas medias hay sistemas reales como Cursor, Claude Code, mobile agent y deep analysis;
  \item Aguas abajo hay escenarios de uso claros: análisis de datos empresariales, ingeniería de software, automatización de GUI.
\end{itemize}
\end{importantbox}

En comparación con las rondas anteriores de auge del «prompt engineering», lo que más le importa a este panel no es «cómo escribir el prompt», sino «cómo convertir el modelo en una capacidad de sistema que pueda trabajar a largo plazo, interactuar con el entorno y encargarse de flujos de trabajo». Este es también el trasfondo por el que se distingue una y otra vez entre «agent» y «agentic».

\begin{knowledgebox}{La diferencia entre Agent y Agentic}
Los ponentes insistieron en que el Agent temprano era, sobre todo, apilar flujos de trabajo sobre un marco de prompting; mientras que Agentic pone más énfasis en sistematizar capacidades como la planeación, la interacción, la recuperación y la exploración, e incluso interiorizarlas más a fondo en los parámetros del modelo. Dicho de otro modo, lo primero se inclina hacia un «andamiaje externo» y lo segundo hacia una «construcción de capacidad nativa».
\end{knowledgebox}

\subsection{El objeto de la discusión pasó de la conversación de un solo turno a la productividad de largo plazo}
El ejemplo que dio Zhang Shaolei (张少磊) es representativo. Su equipo desplazó la atención de un modo de chat de pregunta y respuesta hacia escenarios de automatización de data science: las empresas tienen grandes necesidades de bases de datos y de análisis; el proceso tradicional requiere que una persona escriba SQL a mano, limpie los datos y verifique los resultados repetidamente, mientras que un sistema agentic puede enlazar en un ciclo cerrado la codificación, el análisis, la evaluación y la extracción de conclusiones. La discusión apuntó a un cambio clave: el Agent ya no solo responde preguntas, sino que realiza en lugar de la persona varios pasos de la cadena de trabajo.

\begin{warningbox}{No hay que equiparar al Agent con un chatbot que llama herramientas}
Si al Agent solo se le entiende como «LLM + llamada de herramientas», es fácil subestimar su complejidad de ingeniería. La parte verdaderamente difícil está en: si la descomposición de la tarea es razonable, si el espacio de herramientas es estable, si la memoria puede conservarse entre turnos, si los errores pueden recuperarse y si la retroalimentación del entorno puede convertirse en señal de entrenamiento.
\end{warningbox}

\subsection{Resumen de la sección}
El punto de partida de este panel es muy claro: 2025 se convirtió en el año del Agent no porque el concepto se haya renovado, sino porque por primera vez la capacidad del modelo, la ingeniería de sistemas y la demanda industrial formaron un ciclo cerrado. Todas las discusiones que siguen giran en torno a qué tramo de ese ciclo cerrado es el más corto y cuál conviene reforzar primero.
